\begin{table}[htbp]
\centering
\caption{New Mexico dispensary-level border effect, noise-minimizing specification (grid-search selected)}
\label{tab:nm_dispensary_border_best}
\footnotesize
\begin{tabular}{lccccc}
\toprule
Outcome & Coef. (\%) & SE & $p$ (cl.) & $n$ & Brand FE \\
\midrule
Revenue per capita & $+0.149$ ($+16.1\%$) & 0.107 & 0.162 & 10725 & Yes \\
Revenue (per dispensary) & $+0.150$ ($+16.2\%$) & 0.103 & 0.144 & 10739 & Yes \\
Local dispensary density & $-0.154$ ($-14.3\%$) & 0.160 & 0.336 & 10739 & No \\
\bottomrule
\end{tabular}
\begin{minipage}{0.98\textwidth}
\vspace{0.3em}
\footnotesize
\textit{Notes:} Specification chosen by a full factorial grid search over population control (county vs.\
25-mile-radius), WLS age-weighting (on/off), chain/brand fixed effects (on/off), and clustering level
(county vs.\ ZIP), ranked by standard error on the border coefficient (48 regressions; see
\_grid\_search\_noise.py). County-level clustering dominates ZIP-level clustering in every one of the 16
combinations tested, for all three outcomes. Brand fixed effects (mechanically extracted from each
dispensary's business name) reduce the standard error for the two revenue outcomes by roughly 15-20 percent
despite only 31 percent of the sample identifying the border coefficient once brand FE is added -- the
variance reduction from removing persistent firm-level heterogeneity dominates the loss of identifying
sample. Brand FE does not help local dispensary density (a spatial-competition measure with no natural
brand determinant) and is therefore excluded for that outcome. Local dispensary density itself uses a
25-mile radius (vs.\ the original 10 miles), which alone cuts its standard error from 0.272 to 0.160 by
eliminating a small-denominator blowup for dispensaries in sparsely populated areas.
\end{minipage}
\end{table}
